\subsection{模型二的建立与求解：稳健时空标定与异步融合}\label{sec:model-q2}

\subsubsection{模型二的建立}\label{sec:q2-establishment}

\noindent\textbf{（1）含噪观测模型与时间映射。}

设机器人在参考时间轴上的真实连续轨迹为
\(\mathbf p^\star(t)=(x^\star(t),y^\star(t))^\mathsf T\)。以定位方式 1
的时钟为参考，两种定位方式的离散观测满足

\begin{equation}
\begin{aligned}
\mathbf z_i^{(1)}
  &=\mathbf p^\star\!\left(t_i^{(1)}\right)
    +\boldsymbol\varepsilon_i^{(1)},\\
\mathbf z_j^{(2)}
  &=\mathbf p^\star\!\left(t_j^{(2)}-\Delta t\right)
    +\mathbf b+\boldsymbol\varepsilon_j^{(2)},
\end{aligned}
\label{eq:q2-observation}
\end{equation}

其中，\(\Delta t\) 为方式 2 相对于方式 1 的固定时间偏差，
\(\mathbf b=(b_x,b_y)^\mathsf T\) 为方式 2 相对于方式 1
的固定空间偏差，随机误差满足

\begin{equation}
\mathbb E\!\left(\boldsymbol\varepsilon^{(m)}\right)=\mathbf 0,\qquad
\operatorname{Cov}\!\left(\boldsymbol\varepsilon^{(m)}\right)=\mathbf R_m,
\qquad m=1,2.
\label{eq:q2-noise-assumption}
\end{equation}

由式\eqref{eq:q2-observation}，方式 2 的设备时刻映射到参考时间轴为

\begin{equation}
t_{2,\mathrm{ref}}=t_2-\Delta t.
\label{eq:q2-time-map}
\end{equation}

分别由两组离散观测构造连续位置函数
\(\widetilde{\mathbf z}_1(t)\) 与
\(\widetilde{\mathbf z}_2(t)\)。对候选时间偏差
\(\delta\in\mathcal D\)，其公共支持区间为

\begin{equation}
\mathcal I(\delta)=
[t_{1,\min},t_{1,\max}]
\cap[t_{2,\min}-\delta,t_{2,\max}-\delta].
\label{eq:q2-overlap}
\end{equation}

为使不同候选值在相同样本集合上评价，取固定公共区间

\begin{equation}
\mathcal I_{\mathrm{fix}}
=\bigcap_{\delta\in\mathcal D}\mathcal I(\delta),
\qquad
\mathcal T_{\mathrm{fix}}=\{t_k\}_{k=1}^{N}
\subset\mathcal I_{\mathrm{fix}}.
\label{eq:q2-fixed-grid}
\end{equation}

\noindent\textbf{（2）时间偏差的粗估计。}

对连续位置函数求导，得到两种定位方式的速度序列
\(\mathbf v_m(t)=\mathrm d\widetilde{\mathbf z}_m(t)/\mathrm dt\)。
在粗网格上定义二维速度特征的归一化互相关

\begin{equation}
C(h)=
\frac{\displaystyle
\sum_{t\in\mathcal T_h}
\left(\mathbf v_1(t)-\overline{\mathbf v}_1\right)^\mathsf T
\left(\mathbf v_2(t+h)-\overline{\mathbf v}_2\right)}
{\displaystyle
\sqrt{\sum_{t\in\mathcal T_h}
\left\|\mathbf v_1(t)-\overline{\mathbf v}_1\right\|_2^2}
\sqrt{\sum_{t\in\mathcal T_h}
\left\|\mathbf v_2(t+h)-\overline{\mathbf v}_2\right\|_2^2}},
\label{eq:q2-ncc}
\end{equation}

其中，\(\mathcal T_h\) 为平移后仍位于两组观测支持域内的粗评价网格。
粗时间偏差取为

\begin{equation}
h_c=\arg\max_{h\in\mathcal H} C(h),
\label{eq:q2-coarse-offset}
\end{equation}

并据此在 \(h_c\) 邻域内构造精搜索域 \(\mathcal D\)。

\noindent\textbf{（3）Huber 剖面标定。}

对任意候选 \(\delta\)，在固定公共网格上构造两设备的位置差

\begin{equation}
\mathbf d_k(\delta)=
\widetilde{\mathbf z}_2(t_k+\delta)
-\widetilde{\mathbf z}_1(t_k).
\label{eq:q2-residual}
\end{equation}

对坐标分量 \(r\in\{x,y\}\)，采用 MAD 给出固定稳健尺度

\begin{equation}
s_r=1.4826\,
\operatorname{median}_k
\left|d_{kr}(\delta_0)
-\operatorname{median}_{\ell}d_{\ell r}(\delta_0)\right|,
\label{eq:q2-robust-scale}
\end{equation}

其中 \(\delta_0\) 为精搜索域中点。尺度在整个精搜索过程中保持不变。
令

\begin{equation}
u_{kr}=\frac{d_{kr}(\delta)-b_r}{s_r},\qquad
\rho_c(u)=
\begin{cases}
\dfrac{u^2}{2},& |u|\le c,\\[2mm]
c|u|-\dfrac{c^2}{2},& |u|>c,
\end{cases}
\label{eq:q2-huber-loss}
\end{equation}

则给定 \(\delta\) 时，空间偏差由逐轴 Huber 估计确定：

\begin{equation}
\widehat b_r(\delta)
=\arg\min_{b_r}
\sum_{k=1}^{N}s_r^2
\rho_c\!\left(
\frac{d_{kr}(\delta)-b_r}{s_r}
\right).
\label{eq:q2-profile-bias}
\end{equation}

采用 IRLS 求解式\eqref{eq:q2-profile-bias}。第 \(\ell\) 次迭代的权重与更新量为

\begin{equation}
w_{kr}^{(\ell)}=
\min\!\left(1,\frac{c}{|u_{kr}^{(\ell)}|}\right),
\qquad
b_r^{(\ell+1)}
=\frac{\sum_k w_{kr}^{(\ell)}d_{kr}(\delta)}
{\sum_k w_{kr}^{(\ell)}}.
\label{eq:q2-irls-bias}
\end{equation}

将 \(\widehat{\mathbf b}(\delta)\) 回代，得到一维剖面目标

\begin{equation}
J_H(\delta)=
\frac{2}{N}\sum_{k=1}^{N}\sum_{r\in\{x,y\}}
s_r^2\rho_c\!\left(
\frac{d_{kr}(\delta)-\widehat b_r(\delta)}{s_r}
\right),
\label{eq:q2-huber-objective}
\end{equation}

从而

\begin{equation}
\widehat{\Delta t}
=\arg\min_{\delta\in\mathcal D}J_H(\delta),
\qquad
\widehat{\mathbf b}
=\widehat{\mathbf b}(\widehat{\Delta t}).
\label{eq:q2-calibration-estimator}
\end{equation}

\noindent\textbf{（4）测量噪声协方差估计。}

局部常加速度运动的三阶差分为零。对等间隔位置观测定义

\begin{equation}
\Delta^3\mathbf z_i^{(m)}
=\mathbf z_{i+3}^{(m)}-3\mathbf z_{i+2}^{(m)}
+3\mathbf z_{i+1}^{(m)}-\mathbf z_i^{(m)}.
\label{eq:q2-third-difference}
\end{equation}

若相邻测量误差独立同分布，则每个坐标分量满足
\(\operatorname{Var}(\Delta^3 z_r)=20\sigma_r^2\)。因此

\begin{equation}
\widehat\sigma_{mr}
=\frac{1.4826}{\sqrt{20}}\,
\operatorname{median}_i
\left|
\Delta^3z_{ir}^{(m)}
-\operatorname{median}_{\ell}\Delta^3z_{\ell r}^{(m)}
\right|,
\qquad
\widehat{\mathbf R}_m
=\operatorname{diag}
(\widehat\sigma_{mx}^2,\widehat\sigma_{my}^2).
\label{eq:q2-measurement-covariance}
\end{equation}

\noindent\textbf{（5）异步常加速度状态空间模型。}

将状态写为

\begin{equation}
\mathbf x(t)=
\bigl(\mathbf p(t)^\mathsf T,
\mathbf v(t)^\mathsf T,
\mathbf a(t)^\mathsf T\bigr)^\mathsf T
=(x,y,v_x,v_y,a_x,a_y)^\mathsf T.
\label{eq:q2-state}
\end{equation}

相邻两个观测事件的时间间隔为 \(\tau_k\) 时，状态转移模型为

\begin{equation}
\mathbf x_k=\mathbf F(\tau_k)\mathbf x_{k-1}+\mathbf w_k,
\qquad
\mathbf F(\tau)=
\begin{bmatrix}
\mathbf I_2&\tau\mathbf I_2&\dfrac{\tau^2}{2}\mathbf I_2\\
\mathbf 0&\mathbf I_2&\tau\mathbf I_2\\
\mathbf 0&\mathbf 0&\mathbf I_2
\end{bmatrix}.
\label{eq:q2-transition}
\end{equation}

采用白噪声 jerk 模型，其离散过程噪声协方差为

\begin{equation}
\mathbf Q(\tau;q)=q
\begin{bmatrix}
\dfrac{\tau^5}{20}\mathbf I_2&
\dfrac{\tau^4}{8}\mathbf I_2&
\dfrac{\tau^3}{6}\mathbf I_2\\
\dfrac{\tau^4}{8}\mathbf I_2&
\dfrac{\tau^3}{3}\mathbf I_2&
\dfrac{\tau^2}{2}\mathbf I_2\\
\dfrac{\tau^3}{6}\mathbf I_2&
\dfrac{\tau^2}{2}\mathbf I_2&
\tau\mathbf I_2
\end{bmatrix},
\label{eq:q2-process-covariance}
\end{equation}

其中 \(q\) 为 jerk 谱密度。令
\(\mathbf H=[\mathbf I_2\ \mathbf 0\ \mathbf 0]\)。
方式 2 的事件时刻按式\eqref{eq:q2-time-map}校正，位置观测减去
\(\widehat{\mathbf b}\)；两路原始事件按校正时刻合并排序，不将插值点作为额外观测。

对第 \(k\) 个事件，Kalman 预测与创新为

\begin{equation}
\begin{aligned}
\widehat{\mathbf x}_{k|k-1}
&=\mathbf F_k\widehat{\mathbf x}_{k-1|k-1},\\
\mathbf P_{k|k-1}
&=\mathbf F_k\mathbf P_{k-1|k-1}\mathbf F_k^\mathsf T+\mathbf Q_k,\\
\boldsymbol\nu_k
&=\mathbf y_k-\mathbf H\widehat{\mathbf x}_{k|k-1},\\
\mathbf S_{0,k}
&=\mathbf H\mathbf P_{k|k-1}\mathbf H^\mathsf T+\mathbf R_{m(k)}.
\end{aligned}
\label{eq:q2-kalman-prediction}
\end{equation}

定义门控前归一化创新平方

\begin{equation}
\gamma_k=\boldsymbol\nu_k^\mathsf T
\mathbf S_{0,k}^{-1}\boldsymbol\nu_k.
\label{eq:q2-nis}
\end{equation}

设门限 \(g_\alpha=\chi^2_{2,\alpha}\)，按

\begin{equation}
\lambda_k=\max\!\left(1,\frac{\gamma_k}{g_\alpha}\right),
\qquad
\widetilde{\mathbf R}_k=\lambda_k\mathbf R_{m(k)}
\label{eq:q2-robust-gate}
\end{equation}

放大异常观测的协方差。随后计算

\begin{equation}
\begin{aligned}
\mathbf S_k
&=\mathbf H\mathbf P_{k|k-1}\mathbf H^\mathsf T
+\widetilde{\mathbf R}_k,\\
\mathbf K_k
&=\mathbf P_{k|k-1}\mathbf H^\mathsf T\mathbf S_k^{-1},\\
\widehat{\mathbf x}_{k|k}
&=\widehat{\mathbf x}_{k|k-1}+\mathbf K_k\boldsymbol\nu_k,\\
\mathbf P_{k|k}
&=(\mathbf I-\mathbf K_k\mathbf H)\mathbf P_{k|k-1}.
\end{aligned}
\label{eq:q2-kalman-update}
\end{equation}

\noindent\textbf{（6）过程噪声选择与 RTS 平滑。}

过程噪声调参阶段关闭式\eqref{eq:q2-robust-gate}中的协方差膨胀。令

\begin{equation}
\mathbf e_k=
\operatorname{chol}(\mathbf S_{0,k})^{-1}\boldsymbol\nu_k,
\label{eq:q2-whitened-innovation}
\end{equation}

并在候选集 \(\mathcal Q\) 上最小化

\begin{equation}
S(q)=
\left|\overline{\gamma}(q)-2\right|
+2\,\overline{\left|\rho_1(\mathbf e;q)\right|},
\qquad
\widehat q=\arg\min_{q\in\mathcal Q}S(q),
\label{eq:q2-process-score}
\end{equation}

使二维 NIS 均值接近自由度 2，同时降低白化创新的一阶相关。
固定 \(\widehat q\) 后运行稳健门控和完整前向滤波。

离线融合采用 RTS 后向平滑。对 \(k=K-1,\ldots,1\)，有

\begin{equation}
\begin{aligned}
\mathbf G_k
&=\mathbf P_{k|k}\mathbf F_{k+1}^\mathsf T
\mathbf P_{k+1|k}^{-1},\\
\widehat{\mathbf x}_{k|K}
&=\widehat{\mathbf x}_{k|k}
+\mathbf G_k
\left(\widehat{\mathbf x}_{k+1|K}
-\widehat{\mathbf x}_{k+1|k}\right),\\
\mathbf P_{k|K}
&=\mathbf P_{k|k}
+\mathbf G_k
\left(\mathbf P_{k+1|K}-\mathbf P_{k+1|k}\right)
\mathbf G_k^\mathsf T.
\end{aligned}
\label{eq:q2-rts}
\end{equation}

最后仅在两路数据的公共支持域内，以
\(\Delta t_{\mathrm{out}}=0.1\,\mathrm s\) 传播平滑状态，生成 10 Hz
的位置、速度与加速度序列。

\subsubsection{模型二的求解}\label{sec:q2-solution}

\begin{pseudocodeblock}{算法2\quad 稳健时空标定与异步融合}
{\heiti 输入：}两路含噪定位序列
\(\mathcal Z_1,\mathcal Z_2\)，粗搜索集 \(\mathcal H\)，
精搜索半宽 \(L\)，Huber 阈值 \(c\)，IRLS 次数 \(M\)，
过程噪声候选集 \(\mathcal Q\)，门控概率 \(\alpha\)。\par
{\heiti 输出：}时间偏差 \(\widehat{\Delta t}\)、相对空间偏差
\(\widehat{\mathbf b}\) 及 10 Hz 融合状态序列。\par
\begin{enumerate}[label={\heiti 步骤~\arabic*:},leftmargin=5.5em,itemsep=1pt,topsep=3pt]
  \item 构造两路连续位置函数及速度特征，按式\eqref{eq:q2-ncc}计算
  \(C(h)\)，由式\eqref{eq:q2-coarse-offset}得到 \(h_c\)；
  \item 令 \(\mathcal D=[h_c-L,h_c+L]\)，按式\eqref{eq:q2-fixed-grid}
  构造固定公共评价网格，并由式\eqref{eq:q2-robust-scale}确定固定尺度；
  \item 对每个 \(\delta\in\mathcal D\)，迭代 \(M\) 次
  式\eqref{eq:q2-irls-bias}，计算剖面目标 \(J_H(\delta)\)；
  \item 连续最小化式\eqref{eq:q2-huber-objective}，按
  式\eqref{eq:q2-calibration-estimator}得到
  \(\widehat{\Delta t}\) 与 \(\widehat{\mathbf b}\)；
  \item 由式\eqref{eq:q2-third-difference}和
  式\eqref{eq:q2-measurement-covariance}估计
  \(\widehat{\mathbf R}_1,\widehat{\mathbf R}_2\)；
  \item 关闭门控反馈，对 \(q\in\mathcal Q\) 运行前向滤波，按
  式\eqref{eq:q2-process-score}选择 \(\widehat q\)；
  \item 按校正时刻合并原始事件，依次执行
  式\eqref{eq:q2-kalman-prediction}--\eqref{eq:q2-kalman-update}；
  \item 执行式\eqref{eq:q2-rts}的后向平滑，在公共支持域内按
  0.1 s 间隔输出融合状态序列。
\end{enumerate}
\end{pseudocodeblock}

噪声刻画、时空联合标定和异步状态估计之间的关系如
图\ref{fig:q2-framework}所示。
\begin{figure}[htbp]
  \centering
  \includegraphics[width=.98\textwidth]{q2/method_flow.pdf}
  \caption{问题二稳健标定与异步融合流程}
  \label{fig:q2-framework}
\end{figure}

\subsubsection{结果分析}\label{sec:q2-results}

\noindent\textbf{（1）时间与空间偏差估计。}

两种定位方式分别包含 3000 个和 3948 个有效观测。归一化互相关得到粗时间
偏差 \(h_c=52.72\,\mathrm s\)，据此取精搜索区间
\([42.72,62.72]\,\mathrm s\)。固定公共评价网格含 6897 个 10 Hz
时刻，Huber 阈值取 \(c=1.345\)，逐轴 IRLS 迭代 4 次。

图\ref{fig:q2-objective}给出了精搜索区间内的目标函数剖面。Huber
剖面目标在区间内部取得唯一最低点，连续精修值为

\begin{equation}
\widehat{\Delta t}=50.171732\ \mathrm s,
\qquad
t_{2,\mathrm{ref}}=t_2-50.171732\ \mathrm s.
\label{eq:q2-calibration-result}
\end{equation}

\begin{figure}[htbp]
  \centering
  \includegraphics[width=.84\textwidth]{q2/objective_profile.pdf}
  \caption{问题二时空标定目标函数剖面}
  \label{fig:q2-objective}
\end{figure}

将式\eqref{eq:q2-calibration-result}代入逐轴 Huber 估计，得到

\begin{equation}
\widehat{\mathbf b}
=(3.466998,-1.833200)^\mathsf T\ \mathrm m,
\qquad
\|\widehat{\mathbf b}\|_2=3.921824\ \mathrm m.
\label{eq:q2-bias-result}
\end{equation}

该估计表明，方式 2 相对于方式 1 在 \(x\) 方向存在正向平移，在 \(y\)
方向存在负向平移。由于缺少外部真值，\(\widehat{\mathbf b}\)
仅表示两设备的相对系统偏差，不用于判定单个设备的绝对误差。

\noindent\textbf{（2）时空校正效果。}

表\ref{tab:q2-correction}列出校正前后两设备位置差的统计量。

固定网格上的 Huber 标定 RMSE 为 \(1.45262\,\mathrm m\)。
在校正后的 10 Hz 重叠网格上重新计算，两设备位置差的二维 RMSE 为
\(1.45918\,\mathrm m\)。两者的评价网格不同，均表示含噪观测之间的差异，
不是融合轨迹相对于真实轨迹的绝对误差。

\begin{table}[htbp]
\centering
\caption{问题二时空校正前后的两设备位置差}
\label{tab:q2-correction}
\begin{tabular}{lcc}
\toprule
统计量 & 校正前 & 校正后\\
\midrule
\(\overline{\Delta x}\) / m & \(3.47216\) & \(0.00516\)\\
\(\overline{\Delta y}\) / m & \(-1.83035\) & \(0.00286\)\\
\bottomrule
\end{tabular}
\end{table}

校正前后的二维位置关系见图\ref{fig:q2-calibration}。校正前两条轨迹存在明显
平移，按式\eqref{eq:q2-calibration-result}和式\eqref{eq:q2-bias-result}
同时校正时间与空间偏差后，两组含噪观测围绕同一运动轨迹分布。

\begin{figure}[htbp]
  \centering
  \includegraphics[width=.90\textwidth]{q2/calibration_comparison.pdf}
  \caption{问题二时空校正前后的两设备二维位置}
  \label{fig:q2-calibration}
\end{figure}

\noindent\textbf{（3）噪声参数与异步融合。}

三阶差分的 MAD 稳健估计给出

\begin{equation}
\widehat{\mathbf R}_1
=\operatorname{diag}(0.68991,0.71228)\ \mathrm{m}^2,
\qquad
\widehat{\mathbf R}_2
=\operatorname{diag}(0.65003,0.77553)\ \mathrm{m}^2.
\label{eq:q2-measurement-noise-result}
\end{equation}

过程噪声候选集取
\(\mathcal Q=[10^{-11},10^2]\) 上的对数网格。式\eqref{eq:q2-process-score}
在

\begin{equation}
\widehat q=3.162\times10^{-8}
\label{eq:q2-process-noise-result}
\end{equation}

处取得最小值，该值位于候选网格内部。固定 \(\widehat q\) 后采用
\(\alpha=0.99\) 的 NIS 门控。门控前 NIS 均值为 1.9647、95\%
分位为 5.7980，门控后 NIS 均值为 1.9481；白化创新一阶相关绝对值均值为
0.0153，共有 0.892\% 的观测触发协方差放大。

\noindent\textbf{（4）10 Hz 融合轨迹。}

异步 Kalman 滤波与 RTS 平滑后的公共输出区间为
\(102.0\sim951.5\,\mathrm s\)。按 0.1 s 间隔共生成 8496 个状态点，
每个状态同时包含位置、速度和加速度，全部分量均为有限值。
图\ref{fig:q2-trajectory}给出了 10 Hz 融合位置轨迹；起点与终点标记分别对应
公共支持区间的首末时刻。

\begin{figure}[H]
  \centering
  \includegraphics[width=.80\textwidth]{q2/fused_trajectory_10hz.pdf}
  \caption{问题二异步 Kalman/RTS 的 10 Hz 融合轨迹}
  \label{fig:q2-trajectory}
\end{figure}
